% ============================================================
%  第 8 章  项目管理
% ============================================================
\section{项目管理}
\label{sec:management}

\subsection{开发流程与进度安排}

项目按课程实践指导书规定的四个阶段推进，实际执行的日程安排见
\cref{fig:gantt}。图中标注了两个关键节点：技术风险验证前置于第三日完成，
核心算法在第六日形成可验证的实现。

\begin{figure}[htbp]
\centering
\begin{tikzpicture}[x=7.6mm, y=-6.2mm]

% ---- 时间轴 ----
\foreach \d in {1,...,14}
  \node[font=\scriptsize, color=aegisgray] at (\d,-0.6) {D\d};
\draw[rulegray!55, line width=0.4pt] (0.5,-0.15) -- (14.5,-0.15);

% ---- 阶段背景 ----
\fill[aegisblue!5]  (0.6,0.3) rectangle (4.4,1.7);
\fill[aegisteal!5]  (4.6,0.3) rectangle (11.4,7.7);
\fill[aegisamber!5] (11.6,0.3) rectangle (14.4,9.7);

% ---- 任务条 ----
\newcommand{\task}[5]{
  \fill[#5, rounded corners=1.2pt] (#1-0.38,#3-0.30) rectangle (#2+0.38,#3+0.30);
  \node[anchor=west, font=\scriptsize] at (14.8,#3) {#4};
}

\task{1}{2}{1}{需求分析与威胁建模}{aegisblue!55}
\task{3}{4}{2}{架构设计与技术验证}{aegisblue!55}
\task{5}{6}{3}{检测引擎内核实现}{aegisteal!58}
\task{7}{7}{4}{语义净化与辅助检测}{aegisteal!58}
\task{8}{8}{5}{受控靶场构建}{aegisteal!58}
\task{9}{9}{6}{RASP 探针实现}{aegisteal!58}
\task{10}{10}{7}{网关与控制台服务}{aegisteal!58}
\task{11}{12}{8}{可视化界面实现}{aegisind!52}
\task{13}{13}{9}{安全测试与缺陷修复}{aegisamber!58}
\task{14}{14}{10}{报告撰写与成果整理}{aegisamber!58}

% ---- 里程碑 ----
\node[star, star points=5, star point ratio=2.2, fill=aegisred!85,
      inner sep=1.1pt] at (4,2) {};
\node[star, star points=5, star point ratio=2.2, fill=aegisred!85,
      inner sep=1.1pt] at (6,3) {};
\node[font=\scriptsize, color=aegisred, anchor=west] at (14.8,10.9)
  {\ding{72} 关键里程碑};

% ---- 阶段标注 ----
\node[font=\scriptsize\bfseries, color=aegisblue] at (2.5,-1.5) {设计阶段};
\node[font=\scriptsize\bfseries, color=aegisteal] at (8.0,-1.5) {实现阶段};
\node[font=\scriptsize\bfseries, color=aegisamber] at (13.0,-1.5) {验证阶段};

\end{tikzpicture}
\caption{项目开发进度甘特图}
\label{fig:gantt}
\end{figure}

\subsection{风险管理}

项目启动时识别了六项风险并预设应对措施，实际发生情况见
\cref{tab:risk-register}。其中字节码注入的技术不确定性被判定为最高风险，
故将其验证前置至第三日，并预设了退化方案。

\begin{table}[htbp]
\centering
\caption{风险登记表与实际发生情况}
\label{tab:risk-register}
\footnotesize
\begin{tabular}{@{}l c c p{3.6cm} l@{}}
\toprule
\textbf{风险项} & \textbf{概率} & \textbf{影响} & \textbf{应对措施} & \textbf{实际情况} \\
\midrule
字节码注入受阻     & 中 & 高 & 退化为数据访问层显式埋点 & 未发生，一次验证通过 \\
探针依赖冲突       & 中 & 高 & 打包时排除日志实现       & \textbf{发生并修复} \\
语法解析方言差异   & 高 & 中 & 容错解析加词法层降级     & 发生，按预案处理 \\
图渲染性能不足     & 低 & 中 & 限制节点数并支持折叠     & 未发生 \\
进度延误           & 中 & 高 & 按优先级裁剪增强功能     & 未发生 \\
误报率超标         & 中 & 中 & 提高相似度容差与接口豁免 & \textbf{发生并修复} \\
\bottomrule
\end{tabular}
\end{table}

\subsection{工时统计}

各阶段的工时投入见 \cref{fig:effort}。实现阶段占比最高，
其中检测引擎内核与可视化界面各占实现工时的三成左右。
测试与缺陷修复的工时占比达 18\%，这一比例反映了安全类项目对验证的高要求。

\begin{figure}[htbp]
\centering
\begin{tikzpicture}
\begin{axis}[
  width=11.0cm, height=5.0cm,
  axis lines=left,
  xbar stacked,
  bar width=17pt,
  xlabel={工时占比 (\%)},
  symbolic y coords={{工时分布}},
  ytick=\empty,
  xmin=0, xmax=100,
  legend style={at={(0.5,-0.42)}, anchor=north, legend columns=3,
                font=\scriptsize},
  xmajorgrids=true,
  axis line style={rulegray!85},
  nodes near coords,
  nodes near coords style={font=\scriptsize, color=white},
  every node near coord/.append style={anchor=center}
]
\addplot[fill=aegisblue!60, draw=aegisblue!85] coordinates {(14,{工时分布})};
\addplot[fill=aegisteal!60, draw=aegisteal!85] coordinates {(38,{工时分布})};
\addplot[fill=aegisind!55, draw=aegisind!80]  coordinates {(20,{工时分布})};
\addplot[fill=aegisamber!60, draw=aegisamber!85] coordinates {(18,{工时分布})};
\addplot[fill=aegisgray!45, draw=aegisgray!75] coordinates {(10,{工时分布})};
\legend{需求与设计, 后端实现, 前端实现, 测试与修复, 文档撰写}
\end{axis}
\end{tikzpicture}
\caption{各阶段工时投入占比}
\label{fig:effort}
\end{figure}

\subsection{小组成员与分工}

项目由四名成员协作完成。分工遵循两条原则：其一，按子系统划分职责，
使每位成员对所负责模块拥有完整的设计、实现与测试所有权，避免多人交叉
修改同一模块引入的协调成本；其二，核心算法与总体架构由组长统一负责，
以保证各模块在接口约定与判定逻辑上的一致性。

各成员的职责范围与主要产出见表 \ref{tab:team-roles}。

\begin{table}[htbp]
\centering
\caption{小组成员分工与主要产出}
\label{tab:team-roles}
\footnotesize
\begin{threeparttable}
\begin{tabular}{@{}l l p{5.6cm} p{3.6cm}@{}}
\toprule
\textbf{成员} & \textbf{角色} & \textbf{职责范围} & \textbf{主要产出} \\
\midrule
黄思凯 & 组长 &
  选题论证与总体方案设计；结构指纹与结构差分算法的设计与实现；
  双层架构设计与 RASP 探针实现；三轮测试方法设计；报告统稿 &
  检测引擎内核、探针模块、算法形式化定义 \\
\midrule
罗锦皓 & 组员 &
  安全网关的设计与实现：反向代理、检测过滤器链、参数判定的工程化；
  令牌桶与滑动窗口限流、来源信誉评估；性能测试 &
  网关模块、限流与信誉组件、性能测试报告 \\
\midrule
曾铭 & 组员 &
  受控靶场的构建：业务功能实现、漏洞与安全双实现对照；
  可视化前端：态势总览、语法树差分渲染、攻防演练台 &
  靶场模块、前端五个功能页面 \\
\midrule
吴靖翔 & 组员 &
  控制台服务：事件汇聚、哈希链审计、统计聚合与实时推送；
  OWASP 测试用例的执行、结果记录与回归测试补充 &
  控制台模块、测试脚本与测试记录 \\
\bottomrule
\end{tabular}
\begin{tablenotes}[flushleft]\footnotesize
\item 各成员共同参与需求评审、代码互审与联调测试；报告由组长统稿，
各成员分别撰写本人所负责模块的初稿。
\end{tablenotes}
\end{threeparttable}
\end{table}

协作流程上，项目采用\zh{每日对齐、集中联调}的工作方式：实现阶段每日
简短对齐各自进度与阻塞点；模块间接口（事件上报格式、追踪标识约定、
策略同步协议）在编码前以文档形式冻结，之后各模块并行开发；
集成阶段集中联调，期间发现并修复的跨模块缺陷均记录于开发日志。

模块负责人对其产出的正确性负第一责任，关键设计决策（如探针钩取点的
选择、哈希链字段的构成）经组内讨论后由组长裁定，保证了决策效率与
方案一致性的平衡。

\subsection{版本控制与过程留痕}

项目全程使用 Git 进行版本管理，遵循约定式提交规范，提交信息按
\code{feat}、\code{fix}、\code{docs}、\code{test}、\code{chore} 分类。
代码托管于公开仓库，包含完整的提交历史、设计文档与测试脚本，
构成可追溯的开发过程记录。

\subsection{质量保障措施}

项目采取三项质量保障措施：其一，检测引擎的每项算法均配套单元测试，
且测试用例直接对应算法的形式化性质（如指纹的三条性质各有对应测试）；
其二，每次发现缺陷后补充相应的回归测试，防止问题重现——
本文的命令注入误报修复即配套新增了 12 项测试；其三，
以三轮测试法进行端到端验证，避免单点验证的片面性。
